% Final Review shared style
% Mathematics for Science Quiz -- 2026 Science War

\usepackage[a4paper,margin=1in,headheight=14pt,headsep=17pt,footskip=16mm]{geometry}
\usepackage{fontspec}
\usepackage{xetexko}
\IfFontExistsTF{Noto Serif CJK KR}{%
  \setmainhangulfont[BoldFont={Noto Serif CJK KR Bold},ItalicFont={Noto Serif CJK KR}]{Noto Serif CJK KR}%
}{%
  \setmainhangulfont[ItalicFont={UnBatang}]{UnBatang}%
}
\IfFontExistsTF{Noto Sans CJK KR}{\setsanshangulfont{Noto Sans CJK KR}}{\setsanshangulfont{UnDotum}}
\IfFontExistsTF{Liberation Serif}{%
  \setmainfont{Liberation Serif}%
}{%
  \IfFontExistsTF{Times New Roman}{\setmainfont{Times New Roman}}{}%
}
\IfFontExistsTF{Liberation Sans}{%
  \setsansfont{Liberation Sans}%
}{%
  \IfFontExistsTF{Arial}{\setsansfont{Arial}}{}%
}

\usepackage{microtype}
\microtypesetup{protrusion=false,expansion=false}
\usepackage{amsmath,amssymb,amsthm,mathtools}
\usepackage{array,booktabs,longtable,tabularx}
\usepackage{enumitem,multicol,ragged2e,xcolor,fancyhdr,tocloft,needspace,etoolbox}
\usepackage[unicode,hidelinks]{hyperref}

\definecolor{POSTECHRed}{RGB}{166,25,85}
\definecolor{KAISTBlue}{RGB}{0,65,145}
\definecolor{RuleGray}{RGB}{155,155,155}
\definecolor{SoftGray}{RGB}{95,95,95}
\newcommand{\institutiontext}[2]{\texorpdfstring{\textcolor{#1}{#2}}{#2}}
\DeclareRobustCommand{\POSTECH}{\institutiontext{POSTECHRed}{POSTECH}}
\DeclareRobustCommand{\KAIST}{\institutiontext{KAISTBlue}{KAIST}}

\hypersetup{
  pdftitle={Final Review - Mathematics for Science Quiz},
  pdfauthor={Jungwoo Kim},
  pdfsubject={2026 KAIST-POSTECH Science War Science Quiz Mathematics Final Review},
  pdfcreator={XeLaTeX}
}

% Canonical notation synchronized with MATH000 / the handbook.
\newcommand{\N}{\mathbb{N}}
\newcommand{\Nzero}{\mathbb{N}_0}
\newcommand{\Z}{\mathbb{Z}}
\newcommand{\Q}{\mathbb{Q}}
\newcommand{\R}{\mathbb{R}}
\newcommand{\C}{\mathbb{C}}
\newcommand{\F}{\mathbb{F}}
\newcommand{\Hquat}{\mathbb{H}}
\newcommand{\Pow}{\mathcal{P}}
\newcommand{\dd}{\mathop{}\!\mathrm{d}}
\newcommand{\Exp}{\mathbb{E}}
\newcommand{\Prob}{\mathbb{P}}
\newcommand{\indicator}[1]{\mathbf{1}_{#1}}
\newcommand{\ones}{\mathbf{1}}
\DeclareMathOperator{\rk}{rk}
\DeclareMathOperator{\tr}{tr}
\DeclareMathOperator{\diag}{diag}
\DeclareMathOperator{\Span}{span}
\DeclareMathOperator{\Null}{Null}
\DeclareMathOperator{\Col}{Col}
\DeclareMathOperator{\spec}{spec}
\DeclareMathOperator{\interior}{int}
\DeclareMathOperator{\Res}{Res}
\DeclareMathOperator{\Arg}{Arg}
\DeclareMathOperator{\Log}{Log}
\DeclareMathOperator{\RePart}{Re}
\DeclareMathOperator{\ImPart}{Im}
\DeclareMathOperator{\Var}{Var}
\DeclareMathOperator{\Cov}{Cov}
\DeclareMathOperator{\lcm}{lcm}
\DeclareMathOperator{\ord}{ord}
\DeclareMathOperator{\GL}{GL}
\DeclareMathOperator{\SL}{SL}
\DeclareMathOperator{\Aut}{Aut}
\DeclareMathOperator{\Gal}{Gal}
\DeclareMathOperator{\Bern}{Bern}
\DeclareMathOperator{\Bin}{Bin}
\DeclareMathOperator{\Geom}{Geom}
\DeclareMathOperator{\Poi}{Pois}
\DeclareMathOperator{\Unif}{Unif}
\DeclareMathOperator{\ExpDist}{Exp}
\newcommand{\Normal}{\mathcal{N}}

\setlength{\parindent}{1.18em}
\setlength{\parskip}{0pt}
\setlength{\emergencystretch}{1.4em}
\linespread{1.06}
\setlist[itemize]{leftmargin=1.5em,itemsep=0.12em,topsep=0.28em,parsep=0pt,partopsep=0pt}
\setlist[enumerate]{leftmargin=1.7em,itemsep=0.12em,topsep=0.28em,parsep=0pt,partopsep=0pt}
\newcolumntype{L}[1]{>{\raggedright\arraybackslash}p{#1}}
\newcolumntype{Y}{>{\raggedright\arraybackslash}X}
\setlength{\LTpre}{0.35em}
\setlength{\LTpost}{0.45em}

% Handbook-derived hierarchy, compressed for the final document.
\newcounter{finalchapter}
\newcounter{finalsectionanchor}
\newcommand{\CurrentChapter}{Final Review}
\newcommand{\finalpart}[2]{%
  \clearpage\phantomsection
  \addcontentsline{toc}{part}{Part #1\space #2}%
  \pdfbookmark[0]{Part #1. #2}{part-#1}%
  \thispagestyle{plain}%
  \vspace*{0.10cm}%
  {\noindent\fontsize{16.5}{20}\selectfont\bfseries Part #1\quad #2\par}%
  \vspace{0.34cm}%
  {\color{RuleGray}\hrule height 0.40pt}%
  \vspace{0.58cm}%
}
\newcommand{\finalchapter}[2]{%
  \Needspace{10\baselineskip}%
  \setcounter{finalchapter}{#1}%
  \gdef\CurrentChapter{#1. #2}%
  \markright{#1. #2}%
  \phantomsection
  \addcontentsline{toc}{section}{#1\space #2}%
  \pdfbookmark[1]{#1. #2}{chapter-#1}%
  \vspace{0.58em}%
  {\noindent\fontsize{10.7}{13}\selectfont\bfseries Chapter #1\par}%
  \vspace{0.02em}%
  {\noindent\fontsize{15.2}{18.2}\selectfont\bfseries #2\par}%
  \vspace{0.18em}%
  {\color{RuleGray}\hrule height 0.30pt}%
  \vspace{0.52em}%
}
\newcommand{\finalsection}[1]{%
  \Needspace{6\baselineskip}%
  \stepcounter{finalsectionanchor}%
  \pdfbookmark[2]{#1}{section-\thefinalchapter-\arabic{finalsectionanchor}}%
  \vspace{0.54em}%
  {\noindent\fontsize{12.4}{15}\selectfont\bfseries #1\par}%
  \vspace{0.14em}%
}
\newcommand{\finalheading}[1]{%
  \Needspace{4\baselineskip}%
  \vspace{0.36em}%
  {\noindent\bfseries #1\par}%
  \vspace{0.06em}%
}

% Theorem-like environments: same visual grammar as the handbook.
% Main environments always break after "Theorem. Title" etc.
\newtheoremstyle{finalplain}
  {4.8pt}{4.8pt}{}{}{\bfseries}{}{\newline}
  {\thmname{#1}\thmnote{. #3}}
\newtheoremstyle{finaldefinition}
  {4.8pt}{4.8pt}{}{}{\bfseries}{}{\newline}
  {\thmname{#1}\thmnote{. #3}}
\newtheoremstyle{finalsameline}
  {3.8pt}{3.8pt}{}{}{\bfseries}{.}{0.50em}
  {\thmname{#1}\thmnote{. #3}}

\theoremstyle{finaldefinition}
\newtheorem*{definition}{Definition}
\newtheorem*{formula}{Formula}
\newtheorem*{fact}{Fact}
\newtheorem*{algorithm}{Algorithm}
\newtheorem*{openproblem}{Open Problem}
\theoremstyle{finalplain}
\newtheorem*{theorem}{Theorem}
\newtheorem*{proposition}{Proposition}
\newtheorem*{lemma}{Lemma}
\newtheorem*{corollary}{Corollary}
\theoremstyle{finalsameline}
\newtheorem*{remark}{Remark}
\newtheorem*{strategy}{Strategy}
\newtheorem*{tactic}{Tactic}

\newcommand{\finalblockspace}{\Needspace{5\baselineskip}}
\BeforeBeginEnvironment{definition}{\finalblockspace}
\BeforeBeginEnvironment{formula}{\finalblockspace}
\BeforeBeginEnvironment{fact}{\finalblockspace}
\BeforeBeginEnvironment{algorithm}{\finalblockspace}
\BeforeBeginEnvironment{openproblem}{\finalblockspace}
\BeforeBeginEnvironment{theorem}{\finalblockspace}
\BeforeBeginEnvironment{proposition}{\finalblockspace}
\BeforeBeginEnvironment{lemma}{\finalblockspace}
\BeforeBeginEnvironment{corollary}{\finalblockspace}
\BeforeBeginEnvironment{remark}{\finalblockspace}
\BeforeBeginEnvironment{strategy}{\finalblockspace}
\BeforeBeginEnvironment{tactic}{\finalblockspace}

% Contents: Parts + Chapters only; Section bookmarks remain in the PDF sidebar.
\setcounter{tocdepth}{1}
\renewcommand{\contentsname}{}
\setlength{\cftbeforepartskip}{0.45em}
\renewcommand{\cftpartfont}{\normalsize\bfseries}
\renewcommand{\cftpartpagefont}{\normalsize\bfseries}
\renewcommand{\cftsecfont}{\small}
\renewcommand{\cftsecpagefont}{\small}
\setlength{\cftbeforesecskip}{0.05em}
\setlength{\cftsecindent}{1.25em}
\setlength{\cftsecnumwidth}{0pt}

\pagestyle{fancy}
\fancyhf{}
\fancyhead[L]{\small\textit{Final Review}}
\fancyhead[R]{\small\textit{Mathematics for Science Quiz}}
\fancyfoot[C]{\small\thepage}
\renewcommand{\headrulewidth}{0.35pt}
\renewcommand{\footrulewidth}{0pt}
\fancypagestyle{plain}{%
  \fancyhf{}
  \fancyhead[L]{\small\textit{Final Review}}
  \fancyhead[R]{\small\textit{Mathematics for Science Quiz}}
  \fancyfoot[C]{\small\thepage}
  \renewcommand{\headrulewidth}{0.35pt}%
}

\newcommand{\FinalOpening}{%
  \pagenumbering{gobble}%
  \pagestyle{empty}%
  \thispagestyle{empty}%
  \vspace*{0.48cm}%
  {\noindent\fontsize{21}{25}\selectfont\bfseries Final Review\par}%
  \vspace{0.06cm}%
  {\noindent\fontsize{12.2}{15}\selectfont\itshape Mathematics for Science Quiz\par}%
  \vspace{0.20cm}%
  {\noindent\small 2026 \KAIST--\POSTECH{} Science War\par}%
  \vspace{0.34cm}%
  {\noindent\small Second Edition. 2026년 9월 17일\par}%
  \vspace{0.38cm}%
  {\noindent\normalsize Jungwoo Kim\par}%
  {\noindent\small Department of Mathematics, \POSTECH\par}%
  {\noindent\small 2026 Science Quiz Mathematics Representative\par}%
  \vspace{0.40cm}%
  \hrule height 0.42pt
  \vspace{0.30cm}%
  {\noindent\large\bfseries Contents\par}%
  \vspace{0.05cm}%
  \begingroup
    \small
    \RaggedRight
    \hyphenpenalty=10000\exhyphenpenalty=10000
    \tableofcontents
  \endgroup
  \vfill
  {\noindent\footnotesize\color{SoftGray}
    \textit{Mathematics for Science Quiz}, 2026 Science War Edition (2026년 9월 15일 기준)을 바탕으로 작성.\par
    PPTX 36문항, Problem sets 90문항, Past Problems 103문항의 전수 점검에서 확인된 핵심 결과와 식별 단서를 반영.}
  \thispagestyle{empty}%
  \clearpage
  \pagestyle{fancy}%
  \pagenumbering{arabic}%
}
